% ============================================================
% discussion.tex
% ============================================================

This chapter synthesizes our empirical results, connects them to the broader literature and policy landscape, and organizes the discussion around five substantive themes: the PSLRA as a time-varying procedural filter, the composition-adjusted sensitivity analysis of the raw PSLRA association, judicial geography as a structural determinant of outcomes, the role of case-level characteristics, and the identification limits that constrain the causal interpretation of all estimates.

\section{Policy Implications of the PSLRA Association}

Our results refine rather than overturn the existing PSLRA
literature. Earlier studies had already linked the statute to
stronger merit-based screening and circuit-level variation
\citep{johnson2007, choi2007, perino2003}. The competing-risks framework clarifies the timing of that
association: the dismissal-side effect is concentrated at the
pleading stage rather than spread across the full observation
window, consistent with the statute's procedural design.

The settlement finding is harder to sustain. Single-outcome
and static approaches treat settlement as the primary endpoint
without accounting for dismissal as a competing pathway
\citep{brochet2014, cox2008}. Once that pathway is modeled
jointly and the settlement hazard is subjected to circuit,
time-trend, and cluster-robust checks, the evidence weakens
materially. The picture is therefore asymmetric: the
early-dismissal filter is stable across specifications, but
the case for settlement suppression is qualified.

The limits of that interpretation deserve equal emphasis.
Administrative data record outcomes, not merit. Whether the
cases screened out early are predominantly low-quality suits
or whether some meritorious claims are also being filtered is
a question these data cannot answer. The results speak to
timing and pathway; they are silent on whether the filtering
is normatively warranted.

\section{Composition-Adjusted Sensitivity}

The IPTW triangulation (Section~\ref{sec:results_iptw})
addresses whether the raw PSLRA association survives after
conditioning on observable case-mix changes. The answer
differs sharply by outcome. For dismissal, the weighted
estimate is slightly larger than the raw estimate, so
observable compositional shifts do not explain away the
dismissal association. For settlement, weighting materially
attenuates the negative association. That attenuation is
substantial in magnitude, but it does not identify a
decomposition. The exercise compares point estimates under
different adjustment strategies; it cannot attribute a fixed
share of the raw settlement gap to compositional shift.

Those adjusted estimates, however, carry three caveats that
limit how much weight they can bear. First, some weighting
covariates may themselves lie on post-PSLRA pathways, so the
adjusted hazard ratio is best read as a conditional
composition-adjusted estimate rather than a total reform
effect. Second, the 1992 placebo test
detects significant pre-existing trends in the pre-PSLRA
period, indicating that the model cannot cleanly separate the
reform from concurrent temporal drift. Third, weighting
sharply reduces the effective size of the pre-PSLRA
comparison group relative to the post-PSLRA cohort, leaving a
fundamental overlap problem that reweighting does not resolve.
These composition-adjusted estimates are therefore best
understood as a heavily trimmed sensitivity check rather than
a well-powered quasi-experimental identification of the
PSLRA's effect.

\section{Circuit Heterogeneity and Judicial Geography}

The dominance of circuit geography in our results---confirmed
by both the cause-specific Cox models and the Fine-Gray
subdistribution analysis---raises fundamental questions
about the role of judicial institutions in shaping litigation
outcomes. The Second Circuit remains a distinctive
low-settlement venue with a high dismissal incidence and a
distinct interaction profile in the PSLRA models, but the
cumulative-incidence plots do not make it the single most
extreme dismissal outlier. The current analysis does not
disaggregate districts within the circuit, so it cannot tell
us how much of this pattern is driven by the Southern
District of New York versus the circuit more broadly.

Several mechanisms might explain this pattern, though none are directly
testable with the current data. One possibility is that the Second Circuit
contains a relatively large share of complex, high-stakes securities cases
and correspondingly sophisticated defense counsel, which could reduce early
settlement probability. Another possibility is that judges with extensive
experience in securities litigation apply demanding standards at the
motion-to-dismiss and summary judgment stages, reducing the expected value
of some cases to plaintiffs. The PSLRA $\times$ Circuit interaction
analysis adds a temporal dimension: within the Second Circuit (the
reference category), the PSLRA is associated with a particularly strong
dismissal-side shift and only a much weaker settlement-side association,
consistent with the possibility that the Second Circuit's already low
settlement baseline leaves less room for additional suppression. Not every
interaction is estimated precisely, however: the Fourth and Ninth Circuit
settlement interactions and the Eleventh Circuit dismissal interaction are
not statistically distinguishable from the Second Circuit baseline.

The strongest dismissal-side attenuations appear outside the
Second Circuit, reinforcing that the PSLRA's practical
association is geographically uneven. Prior studies have
documented circuit-level variation in PSLRA application
\citep{perino2003, choi2007}; our competing-risks framework adds precision about how this
variation manifests jointly across settlement and dismissal
timing, rather than treating each outcome in isolation.

The shared frailty models (Section~\ref{sec:frailty}) add an
important dimension to this geographic analysis. They indicate
modest residual settlement heterogeneity but much less
residual dismissal heterogeneity after the fixed-effect
covariates are included, and the cluster-robust results
reinforce that asymmetry by widening uncertainty for
settlement more than for dismissal.

If the PSLRA's filtering effects vary substantially across
circuits, plaintiffs filing in different jurisdictions face
qualitatively different litigation environments. This raises
questions about regulatory forum shopping, the appropriate
geographic scope of securities law reform, and whether
post-PSLRA federal litigation produces jurisdiction-specific
differences in recovery patterns in practice.

\section{Structural Case Characteristics: MDL, Origin,
and Statutory Basis}

The MDL finding (Section~\ref{sec:results_characteristics})
illustrates why the choice of estimand matters. MDL
consolidation is associated with lower cause-specific
settlement and dismissal hazards, especially dismissal, yet
raises the cumulative probability of settlement in the
Fine-Gray model. The settlement-side Fine-Gray model has weak
held-out discrimination and rejects proportional
subdistribution hazards, so the SHR is best read as a
cumulative-incidence summary rather than a stable predictive
relation. Taken together, the estimates suggest MDL cases are
slower to resolve but substantially less likely to exit
through dismissal.

Cases in the ``Other'' origin category (which includes
removed and reinstated cases) are more settlement-prone and
less dismissal-prone, consistent with the interpretation that
cases entering federal court through non-original pathways
involve claims where settlement is a more natural resolution.
Section~11 cases also appear somewhat more dismissal-prone
than Section~10(b) cases, which may reflect tracing and
standing requirements in IPO and secondary offering contexts
that create additional dismissal opportunities.

\section{Limitations}
\label{sec:limitations}

Several limitations constrain the interpretation of our findings. We address them in order of their bearing on the central
estimates.

\paragraph{Identification limits.}
The key identification issue is whether the rise in dismissals after 1995 reflects the effect of the PSLRA or simply a continuation of pre-existing trends. The 1992 placebo test helps address this concern: when a false cutoff is imposed within the pre-reform period, it produces a significant downward trend in dismissals—opposite in sign to the estimated post-reform effect. This makes it less likely that a simple monotonic pre-trend is driving the main result.

However, this does not fully resolve the issue. A V-shaped secular pattern in dismissal risk remains possible, and the model’s sensitivity to a period with no actual policy change weakens the case for a causal interpretation. The narrow-window estimates point in the same direction as a reform effect, but they do not meet the standards of a true regression discontinuity design. As a result, the findings should be interpreted as associational or composition-adjusted rather than causal.


\paragraph{Post-treatment bias in composition adjustment.}
Several covariates included in the IPTW propensity score
model (circuit filing location, case origin, MDL
consolidation, and statutory basis) may themselves be
consequences of the PSLRA rather than pre-treatment
confounders. If the reform altered where and how securities
class actions were filed, the composition-adjusted estimates
condition on some of the reform's potential downstream
channels. The composition-adjusted HR should therefore be
interpreted as an association conditional on those observed
case-mix shifts, not as the
``true'' reform effect and not as a literal lower bound on
the total PSLRA association.

\paragraph{Unobserved heterogeneity.}
Three categories of unobserved confounders most directly affect the PSLRA interpretation in this analysis. First, the IDB contains no case-merit signals: stock price drops at filing, parallel SEC enforcement actions, accounting restatements, and insider trading allegations are strong predictors of outcomes \citep{cox2008, johnson2007, choi2007} but are absent from our data. The PSLRA's early dismissal association therefore cannot be decomposed into meritless cases filtered out versus meritorious claims caught in the same net. Second, judge-level heterogeneity is entirely unmodeled: the IDB lacks consistent judge identifiers, confining our random-effects analysis to the circuit level. The settlement frailty estimate likely conflates local judicial culture with individual assignment effects that circuit frailties cannot separate. Third, the PSLRA is perfectly collinear with calendar time, meaning secular trends in SEC enforcement intensity, plaintiff-bar sophistication, and broad economic conditions potentially confound the PSLRA association even after flexible spline time-trend controls. These are not generic limitations applicable to any administrative dataset; they are the specific boundary conditions of what the FJC's IDB can credibly identify in the securities litigation context.

\paragraph{Small-cluster frailty caveat.}
The shared frailty models operate with only 11 circuit-level
clusters, well below the 20--30 clusters typically
recommended for reliable variance estimation. The frailty
variance estimates should be treated as indicative rather than
precise, and may be biased downward.

\paragraph{Disposition coding and judgment-bearing disaggregation.}
Among the judgment bearing disposition codes~\{4, 6, 15, 17,
19, 20\}, 93.4\% of the Code~6 plaintiff-victory
reclassifications are post-PSLRA filings, and the broader
six-code disaggregation likely shares this imbalance. Two
explanations are possible. The asymmetry may reflect a genuine
behavioral shift in which courts and parties increasingly used
judgment-bearing codes to record negotiated resolutions,
understating the true post-PSLRA settlement rate across all
three coding schemes. Alternatively, it may reflect a recording artifact stemming
from systematic improvement in IDB \texttt{JUDGMENT}-field
coverage after 1995. Distinguishing these
requires merging IDB records with CourtListener docket text
and comparing match rates across eras---something beyond
the scope of this thesis but with direct implications for the
credibility of all three coding schemes and the magnitude of
the PSLRA settlement association across specifications.

\paragraph{Proportional hazards violations.}
Global Grambsch-Therneau tests reject proportional hazards
in both extended models
(Section~\ref{subsec:ph_diagnostics}). All constant hazard
ratio estimates in this thesis must therefore be interpreted
as time-averaged summaries of a more complex temporal
relationship, not as constant instantaneous risks. The
piecewise Cox model partially addresses this for the PSLRA
covariate, but time-varying effects of other covariates remain
unmodeled. The Fine-Gray models also reject proportional
subdistribution hazards for the PSLRA term and globally, so
the reported SHRs carry the same time-averaged interpretation
as the cause-specific hazard ratios rather than a literally
constant effect over follow-up.

Taken together, the five substantive themes of this chapter
establish a picture of securities litigation that is richer
and more heterogeneous than the standard single-outcome view.
The competing-risks framework materially improves description
and interpretation, but the placebo result, unobserved merit
signals, and settlement fragility keep the final claims on a
strictly associational footing. What remains is to consider
what those claims imply---for litigation research, for
reform evaluation, and for the methodological choices
available to future empirical work.
